\documentclass[11pt,a4paper,openany]{book}

\usepackage[french]{course}

% Deux boîtes propres à ce document, comme dans elec/course.tex :
% goatbox, l'idée d'une section dite le plus simplement possible, et recapbox, le grand
% récapitulatif de fin de chapitre à relire la veille de l'examen (sécable : il est long).
\definecolor{goatColor}{HTML}{8A5A1F}
\definecolor{recapColor}{HTML}{5B3A8E}
\newtcolorbox{goatbox}[1][]{
  enhanced, colback=goatColor!6, colframe=goatColor,
  boxrule=0.5pt, arc=3pt, left=5pt, right=5pt, top=4pt, bottom=4pt,
  fonttitle=\bfseries\small, coltitle=white,
  colbacktitle=goatColor,
  attach boxed title to top left={yshift=-3pt, xshift=6pt},
  boxed title style={size=small, colback=goatColor, boxrule=0pt, arc=2pt},
  title=Expliqué à une chèvre, #1
}
\newtcolorbox{recapbox}[1][]{
  enhanced, breakable, colback=recapColor!5, colframe=recapColor,
  boxrule=1.2pt, arc=4pt, left=8pt, right=8pt, top=8pt, bottom=8pt,
  fonttitle=\bfseries\large, coltitle=white,
  colbacktitle=recapColor,
  title=Récap pour l'examen, #1
}

% Capture d'une diapositive du cours : \slide[largeur]{cNN-PP}{légende}
\graphicspath{{images/course/}}
\newcommand{\slide}[3][0.8]{%
  \begin{center}
    \includegraphics[width=#1\linewidth]{#2}\\[2pt]
    {\small\itshape #3}
  \end{center}}

\begin{document}
\raggedbottom

\coursetitle{Comportement mécanique des matériaux --- expliqué}{BA5}{2026--2027}

\tableofcontents
\clearpage

\chapter*{Comment lire ces notes}
\addcontentsline{toc}{chapter}{Comment lire ces notes}

Ces notes reprennent le cours \emph{MSE 234 --- Comportement mécanique des matériaux}
(A.~Mortensen, EPFL, automne 2026), chapitre par chapitre, à partir des diapositives. Elles ne
recopient pas les diapositives : elles \textbf{expliquent} ce qu'il faut en comprendre, pourquoi
les formules ont cette forme et comment s'en servir. Les exemples chiffrés qui ne viennent pas
des diapositives sont signalés comme tels.

Chaque chapitre contient :
\begin{itemize}
  \item les définitions, lois et formules du cours dans des boîtes de couleur (vert :
        définition, rouge : loi/théorème, bleu : propriété, orange : attention, gris :
        remarque, exemples en encadré) ;
  \item \textbf{les captures des diapositives} utiles (courbes, diagrammes, schémas), avec ce
        qu'il faut y regarder ;
  \item \textbf{des boîtes brunes \og Expliqué à une chèvre\fg{}} : l'idée de la section avec
        les mots les plus simples possibles, à lire d'abord quand une notion ne passe pas ;
  \item \textbf{une grande boîte violette \og Récap pour l'examen\fg{}} à la fin : ce qu'il
        faut savoir, assez court pour être relu la veille, et comment le chapitre s'appuie
        sur les précédents.
\end{itemize}
Le dernier chapitre est un \textbf{récapitulatif général} : les formules à connaître par c\oe
ur, les méthodes pas à pas et les pièges classiques.

\begin{rembox}[title={L'examen en deux mots (chapitre 0 du cours)}]
\begin{itemize}
  \item Examen final \textbf{sans document ni formulaire} (seulement un dictionnaire de
        traduction non technique) : la liste des \og formules à connaître par c\oe ur\fg{}
        est reprise au dernier chapitre.
  \item Note : 75\,\% questions sur le cours, 18\,\% questions sur les TP, 7\,\% cahier de
        laboratoire et participation aux TP.
  \item Ce qui fait foi : le contenu des diapositives. Les livres (Ashby \& Jones, le
        polycopié \emph{Métaux et Alliages}, Callister, Brinson \& Brinson) aident à
        comprendre mais vont au-delà.
\end{itemize}
\end{rembox}

% =========================================================
\chapter{Les essais mécaniques}
% =========================================================

\begin{rembox}[title=En bref]
\textbf{De quoi parle ce chapitre ?} Avant de parler des matériaux, il faut savoir
\textbf{comment on mesure} leur comportement mécanique. Le chapitre passe en revue les essais
classiques : traction (le roi), compression, dureté, flexion, essais où le temps compte
(vitesse de déformation, relaxation, fluage), puis les essais de rupture (ténacité, Charpy,
formabilité, fatigue). \
\textbf{Les trois idées à retenir} : la différence entre grandeurs \emph{nominales} et
\emph{vraies} ; la \textbf{striction} comme instabilité (critère de Considère) ; la
\textbf{mécanique de la rupture} ($K = Y\sigma\sqrt{\pi a}$), qui reviendra dans tous les
chapitres suivants.
\end{rembox}

\section{Pourquoi l'essai de traction est le roi des essais}

\slide[0.75]{c01-06}{Le barreau de traction : une section réduite (\emph{gauge section}) entre deux zones d'amarrage.}

On tire sur un barreau dont le milieu est aminci. Tout l'intérêt est là : dans la section
réduite, l'état de contrainte est le plus simple possible --- \textbf{une seule composante,
$\sigma_{xx}$, la même partout}. Toutes les autres composantes du tenseur des contraintes sont
nulles. On sait donc exactement ce qu'on impose au matériau, ce qui n'est vrai pour presque
aucun autre essai. Les zones d'amarrage sont plus larges pour que ce soit la section réduite,
et elle seule, qui se déforme et casse.

La déformation, elle, a \textbf{trois} composantes : le barreau s'allonge selon $x$ mais
rétrécit selon $y$ et $z$.

\begin{thmbox}[title=Élasticité linéaire isotrope en traction]
\[ \varepsilon_{xx} = \frac{\sigma_{xx}}{E}, \qquad
   \varepsilon_{yy} = \varepsilon_{zz} = -\nu\,\frac{\sigma_{xx}}{E} \]
$E$ = module d'Young (rigidité), $\nu$ = coefficient de Poisson (combien le matériau se
rétrécit latéralement quand on l'allonge).
\end{thmbox}

\begin{propbox}[title={À volume constant, $\nu = 1/2$}]
Le changement relatif de volume vaut, pour de petites déformations,
$\dfrac{\Delta V}{V} \approx \varepsilon_{xx} + \varepsilon_{yy} + \varepsilon_{zz}$.
Si le volume ne change pas (c'est \textbf{grosso modo le cas en déformation plastique}),
cette somme est nulle, donc $\varepsilon_{yy} = \varepsilon_{zz} = -\varepsilon_{xx}/2$ :
tout se passe comme si $\nu = 0{,}5$. En élasticité, $\nu$ est plus petit (souvent
$\approx 0{,}3$ pour les métaux, $\approx 0{,}25$ pour le verre) : le volume augmente un peu
quand on tire.
\end{propbox}

\begin{rembox}[title=En pratique]
La force $P$ est mesurée par une \textbf{cellule de charge}. L'allongement $\Delta L$ se lit
sur le déplacement de la traverse (peu précis, car toute la machine se déforme aussi), ou mieux
avec un \textbf{extensomètre} pincé sur la section réduite, une \textbf{jauge résistive} collée
dessus, ou par suivi optique de points à la surface.
\end{rembox}

\begin{goatbox}
Tirer sur un barreau, c'est la seule façon simple d'être sûr de ce que l'on fait subir au
matériau : partout au milieu du barreau, il est tiré de la même façon, dans une seule
direction. Quand il s'allonge, il maigrit sur les côtés, comme un chewing-gum qu'on étire.
\end{goatbox}

% =========================================================
\chapter{Les céramiques}
% =========================================================

\begin{rembox}[title=En bref]
\textbf{De quoi parle ce chapitre ?} Première grande classe de matériaux : les céramiques
(pierres, briques et porcelaines, verres, ciments et bétons, céramiques techniques). On voit
leurs propriétés (rigides, dures, réfractaires\ldots{} mais \textbf{fragiles}), leur structure
atomique (cristaux ioniques ou covalents, verres amorphes), leur élaboration (argiles,
frittage, mise en forme du verre, contraintes thermiques) et surtout leur
\textbf{comportement mécanique} : rupture fragile, \textbf{statistique de Weibull}, rupture
en compression, fissuration lente. \\
\textbf{Le fil rouge} : tout découle de $K = Y\sigma\sqrt{\pi a}$ (chapitre~1) appliqué à des
matériaux sans plasticité et pleins de petits défauts.
\end{rembox}

\section{Panorama : une famille rigide, dure\ldots{} et fragile}

\begin{defnbox}[title=Les grandes familles de céramiques]
\begin{enumerate}
  \item \textbf{Naturelles} : calcaire/marbre, grès, granite, glace.
  \item \textbf{Poreuses et vitrifiées} : porcelaine, faïence, briques (à base d'argiles).
  \item \textbf{Verres} : verre sodocalcique ($\approx$ 70\,\% SiO$_2$, 10\,\% CaO, 15\,\%
        Na$_2$O), borosilicaté (Pyrex).
  \item \textbf{Ciments et bétons} : ciment Portland (CaO + SiO$_2$ + Al$_2$O$_3$).
  \item \textbf{Céramiques techniques} (haute performance) : alumine Al$_2$O$_3$, SiC,
        Si$_3$N$_4$, zircone ZrO$_2$\ldots{} plus fibres de renfort, revêtements (diamant),
        abrasifs.
\end{enumerate}
\end{defnbox}

Elles sont faites des éléments les plus \textbf{abondants} de l'écorce terrestre (O, Si, Al,
Ca, Na, Mg, K\ldots) : beaucoup sont bon marché.

\subsection*{Choisir un matériau léger et rigide : les indices $E/\rho$ et $E/\rho^2$}

Les cartes d'Ashby placent chaque matériau dans un plan (densité, module\ldots) en échelle
logarithmique. Pour s'en servir, il faut savoir \textbf{quelle combinaison de propriétés} on
veut maximiser, et cela dépend de la forme de la pièce.

\begin{itemize}
  \item \textbf{Barreau en traction} (rayon $r$, longueur $L$, flèche $\delta$ sous la force $F$) :
        $\dfrac{F}{\pi r^2} = E\dfrac{\delta}{L}$ et $M = \pi r^2 L\rho$, d'où
        $M = \dfrac{L^2 F}{\delta}\,\dfrac{\rho}{E}$. Pour la masse la plus faible à rigidité
        donnée : \textbf{maximiser $E/\rho$}.
  \item \textbf{Poutre pleine en flexion} (section carrée $t\times t$, longueur $l$) :
        $\delta = \dfrac{4l^3F}{Et^4}$ et $M = lt^2\rho$, d'où
        $M = \Big(\dfrac{4l^5F}{\delta}\Big)^{1/2}\dfrac{\rho}{E^{1/2}}$ :
        \textbf{maximiser $E^{1/2}/\rho$}, c'est-à-dire $E/\rho^2$.
  \item \textbf{Tube mince en flexion} (rayon $r$ fixé, épaisseur $t \ll r$) : seule $t$ est
        libre, $M \propto \rho/E$ : de nouveau \textbf{$E/\rho$}.
\end{itemize}

Pourquoi la flexion d'une poutre pleine donne-t-elle $\rho^2$ ? Parce qu'en flexion la rigidité
croît comme $t^4$ alors que la masse croît comme $t^2$ : épaissir une poutre légère est très
efficace, donc la densité \og pèse\fg{} deux fois plus.

\slide[0.6]{c02-21}{Carte module--densité : les droites $E/\rho = $ cste et $E/\rho^2 = $ cste. Les céramiques techniques sont en haut : très rigides pour leur poids.}

\begin{propbox}[title=Ce que disent les cartes]
\begin{itemize}
  \item \textbf{Module} : très élevé (alumine 380\,GPa, SiC 410\,GPa, diamant 1050\,GPa).
  \item \textbf{Densité} moyenne (2,2 à 4, sauf zircone et carbures lourds) : excellentes en
        $E/\rho$.
  \item \textbf{Limite d'élasticité} énorme, mais \textbf{sans importance} : elles cassent
        avant de l'atteindre.
  \item \textbf{Réfractaires} (points de fusion élevés) et souvent résistantes à l'attaque
        chimique.
  \item \textbf{Ténacité très faible} ($K_c \approx 0{,}7$ à $\approx 10\,\text{MPa}\sqrt{\text m}$,
        contre 20 à 200 pour les métaux) : \textbf{c'est leur grande faiblesse}.
\end{itemize}
\end{propbox}

\slide[0.6]{c02-26}{Carte ténacité--module : les céramiques sont à droite (rigides) mais en bas (fragiles).}

\begin{goatbox}
Les céramiques, c'est la pierre et le verre : très rigides, très dures, légères pour leur
rigidité, elles supportent la chaleur et la rouille ne les touche pas. Mais elles ont un gros
défaut : la moindre fissure les fait casser net. Pour choisir un matériau, on ne regarde pas
une seule propriété mais un mélange (rigidité divisée par poids, par exemple), et ce mélange
dépend de la forme de la pièce.
\end{goatbox}

% =========================================================
\chapter{Métaux et alliages 1 : l'aluminium, le cuivre et leurs alliages}
% =========================================================

\begin{rembox}[title=En bref]
\textbf{De quoi parle ce chapitre ?} Les métaux, à travers deux exemples : l'\textbf{aluminium}
et le \textbf{cuivre}. On voit d'où ils viennent (ressources, énergie, recyclage, extraction
guidée par le diagramme d'Ellingham), puis le c\oe ur du chapitre : comment la
\textbf{microstructure} se forme (solidification, ségrégation, maturation, précipitation) et
comment on l'utilise pour \textbf{durcir} un métal. \\
\textbf{Les quatre mécanismes de durcissement vus ici} : écrouissage (défait par le recuit),
solution solide, durcissement structural (précipités), renfort par une seconde phase dure. Tous
(sauf le dernier) consistent à mettre des \textbf{obstacles sur le chemin des dislocations}.
\end{rembox}

\section{Panorama des métaux : propriétés, ressources, énergie}

\begin{propbox}[title=Ce qui distingue les métaux et alliages]
Un élément principal (Fe, Al, Cu, Mg, Ti, Ni\ldots) plus des \textbf{éléments d'alliage}. Moins
rigides et moins durs, souvent plus lourds que les céramiques, mais \textbf{ductiles} et
\textbf{tenaces} : c'est ce qui en fait les matériaux de structure par excellence. Ils peuvent
s'oxyder, ont une empreinte énergie/carbone non négligeable, mais se recyclent facilement.
\end{propbox}

\begin{rembox}[title={Une coïncidence utile}]
Fe, Al, Ti, Mg et Ni ont presque le \textbf{même $E/\rho$} ($\approx 25$\,GPa/(Mg/m$^3$)) : en
traction, à rigidité égale, une pièce en aluminium pèse autant qu'une pièce en acier. L'avantage
des alliages légers n'apparaît qu'en flexion ($E/\rho^2$, chapitre~2) ou en résistance.
\end{rembox}

\begin{defnbox}[title={Réserves et ressources}]
\textbf{Réserves} : gisements connus exploitables \textbf{avec profit}, aujourd'hui, avec les
techniques d'aujourd'hui (elles augmentent quand le prix monte ou que la technique progresse).
\textbf{Ressources} : tout ce dont on peut envisager l'extraction, l'autre extrême.
\end{defnbox}

Le coût énergétique d'un métal augmente avec sa \textbf{dilution} dans l'écorce terrestre et
avec la stabilité de son oxyde : l'aluminium, très abondant mais très lié à l'oxygène, coûte
cher en énergie ; l'or, peu lié mais très dilué, encore plus.

\subsection*{Le recyclage est limité par la croissance de la demande}

\slide[0.6]{c03a-38}{Ce qu'on recycle aujourd'hui a été produit il y a $n$ années, quand on produisait moins.}

\begin{thmbox}[title=Taux de recyclage maximal]
Si la production croît de $i$ par an et qu'un objet dure $n$ ans en moyenne, le métal récupérable
aujourd'hui vient de la production d'il y a $n$ ans ($A$), plus faible que celle d'aujourd'hui
($B$) :
\[ F \le \frac{A}{B} = (1+i)^{-n} \]
\end{thmbox}

\begin{exbox}[title={Exemple chiffré (pas dans les diapositives)}]
Aluminium : $n = 15$ ans, croissance $i = 4\,\%$ par an :
$F \le 1{,}04^{-15} = 0{,}56$. Même en recyclant \emph{tout}, on ne couvrirait que 56\,\% de la
production : le reste doit venir du minerai.
\end{exbox}

\begin{goatbox}
Les métaux sont les matériaux qui plient au lieu de casser. On en a plein sous nos pieds, mais
les extraire coûte de l'énergie, surtout quand le métal est très attaché à l'oxygène (aluminium)
ou très dispersé (or). Recycler aide beaucoup, mais tant qu'on en veut chaque année plus, les
vieux objets à recycler ne suffisent pas à tout fournir.
\end{goatbox}

\section{Produire l'aluminium et le cuivre}

\subsection*{Le diagramme d'Ellingham}

\slide[0.7]{c03a-45}{Diagramme d'Ellingham : plus la droite d'un oxyde est basse, plus l'oxyde est stable. Al$_2$O$_3$ est tout en bas.}

\begin{defnbox}[title={Lire le diagramme d'Ellingham}]
Pour chaque réaction \og métal + O$_2$ $\to$ oxyde\fg{} (rapportée à une mole d'O$_2$), on trace
$\Delta G^\circ$ en fonction de $T$.
\begin{itemize}
  \item \textbf{Plus la droite est basse} ($\Delta G^\circ$ très négatif), \textbf{plus
        l'oxyde est stable}, plus il est difficile d'en extraire le métal.
  \item Un élément dont la droite est \textbf{en dessous} de celle d'un autre peut
        \textbf{réduire} l'oxyde de cet autre (il lui \og vole\fg{} l'oxygène).
  \item La droite C + O$_2$ $\to$ 2CO \textbf{descend} quand $T$ augmente : à haute
        température, le carbone devient capable de réduire de nombreux oxydes (fer, cuivre),
        mais pas Al$_2$O$_3$ dans des conditions raisonnables.
\end{itemize}
(Une simplification : les droites dépendent aussi des pressions partielles des gaz.)
\end{defnbox}

\subsection*{Le cuivre : facile à réduire, difficile à trouver}

Le cuivre ne représente que 50\,ppm de l'écorce : le minerai (sulfures, souvent
CuFeS$_2$) est pauvre et coûteux à concentrer. En revanche, Cu est presque \textbf{noble} :
réduire ses composés demande peu d'énergie. La voie \textbf{pyrométallurgique} : concentration
par flottation, fusion en \textbf{matte} (Cu$_2$S--FeS), \textbf{convertisseur} où l'on souffle de
l'air (le fer part dans la scorie avec la silice, le soufre en SO$_2$), cuivre blister, puis
\textbf{raffinage électrolytique} (anodes impures, cathodes de cuivre pur à 99,99\,\%). Il existe
aussi une voie \textbf{hydrométallurgique} (lixiviation à l'acide sulfurique, puis
électrolyse).

\subsection*{L'aluminium : abondant mais très lié à l'oxygène}

Son oxyde est si stable qu'on ne peut pas le réduire au carbone : il faut
l'\textbf{électrolyse}.
\begin{enumerate}
  \item \textbf{Procédé Bayer} : de la bauxite on tire de l'alumine Al$_2$O$_3$ pure.
  \item \textbf{Procédé Hall--Héroult} (1886) : l'alumine est dissoute dans de la cryolithe
        fondue ($\approx 960$\,°C) et électrolysée : $2\,\text{Al}^{3+} + 6e^- \to 2\,\text{Al}$
        à la cathode ; à l'anode en carbone, $3\,\text{O}^{2-} \to \frac32\text{O}_2 + 6e^-$,
        l'oxygène brûlant l'anode en CO$_2$.
\end{enumerate}

\slide[0.6]{c03a-58}{La cuve d'électrolyse à anodes précuites.}

\begin{propbox}[title={L'énergie de l'aluminium}]
Production mondiale $\approx 100$\,Mt/an, dont $\approx 30\,\%$ recyclé. Produire 1\,kg depuis le
minerai coûte $12 \pm 2$\,kWh (limite théorique 6,5\,kWh/kg) ; le recycler, $\approx 1$\,kWh/kg
(10\,\%).
\end{propbox}

\begin{exbox}[title=Le petit exercice des diapositives]
À combien d'accélérations de 0 à 100\,km/h correspondent 12\,kWh/kg ? Par kilogramme,
$\frac12 v^2 = \frac12(10^5/3600)^2 = 386$\,J ; $12$\,kWh $= 43\,200$\,kJ ; donc
$43\,200\,000/386 \approx 112\,000$ accélérations, soit 31 par jour pendant 10 ans. Alléger un
véhicule avec de l'aluminium peut donc rembourser l'énergie de sa production --- à évaluer au
cas par cas.
\end{exbox}

\begin{propbox}[title={Pourquoi l'aluminium ne rouille pas}]
Son oxyde, très stable, forme une couche dure, inerte et \textbf{mécaniquement compatible} avec
le métal : le rapport de \textbf{Pilling--Bedworth} (volume de l'oxyde / volume du métal
consommé) est proche de 1, donc la couche se forme sans fissures ni décollement et protège le
métal. On dit que l'aluminium est \textbf{passivé}.
\end{propbox}

\begin{rembox}[title=Mise en forme]
Fusion, puis soit coulée en lingot suivie de \textbf{laminage} (tôles, feuilles),
\textbf{extrusion} (profilés, tubes) ou \textbf{forgeage} ; soit \textbf{moulage} direct
(sable, cire perdue, moule permanent, sous pression) ; et désormais l'impression 3D (fusion
laser de poudre).
\end{rembox}

\begin{goatbox}
Le diagramme d'Ellingham est un classement de \og qui tient le plus fort à l'oxygène\fg{}.
L'aluminium y tient tellement qu'aucun charbon ne peut le lui arracher : il faut l'électricité,
et beaucoup. Le cuivre, lui, lâche son oxygène (ou son soufre) facilement, mais il est rare et
dispersé. La même passion de l'aluminium pour l'oxygène le protège : une fine peau d'oxyde,
bien collée, l'empêche de rouiller.
\end{goatbox}

% =========================================================
\chapter{Le fer et ses alliages}
% =========================================================

% =========================================================
\chapter{Récapitulatif général pour l'examen}
% =========================================================


\end{document}
